% !TeX root = ../教材题库_全部条目_转写工作稿.tex
% Source: XJ-B1-C03-EX-P081-Y1313
% 第3章 / 3.2 函数的基本性质 / 教材 p.81、82 / PDF 第88、89页
% OCR未校对：公式与文字均可能有误；未完成数学TeX转写；禁止直接发给学生
% 原题及相邻原解截图（核对用）：../题目原图/XJ-B1-C03-EX-P081-Y1313-01.webp, ../题目原图/XJ-B1-C03-EX-P081-Y1313-02.webp
% 本文件只在转写完成后作为题干使用；OCR 条目尚不是可靠题干。
\begin{bankitem}{3.2 例题 例题2}{第3章 / 3.2 函数的基本性质 / 教材 p.81、82 / PDF 第88、89页}{XJ-B1-C03-EX-P081-Y1313}{OCR未校对：公式与文字均可能有误；未完成数学TeX转写；禁止直接发给学生}
\begin{ocrdraft}
[教材p.81 / PDF第88页]
\par
0二0]证明函数丆0）一x十一0>0）在区间（0，1]上单调递减，在区间0，
\par
十亠）上单调递增，并指出函数在区间（0，十。。）上的最值点和最值．
\par
解．设和是区间（0，0上任意两个实数，且厶<，则
\par
1
\par
1工2
\par
1
\par
由0<厶<工2<1，得0<厶<1，
\par
[教材p.82 / PDF第89页]
\par
1
\par
于是
\par
工1主2
\par
因此，函数丆0）在区间（0，1]上单调递减．
\par
O设和是区间0，十亠）上任意两个实数，
\par
且的<2．
\par
1
\par
由2>厶1，得的>1，
\par
1工2
\par
1
\par
于是
\par
工1主2
\par
因此，函数不到在区间[1，十亠）上单调递增．
\par
综合．O可知，在区间（0，十。。）上，函数不x）在x=l处取到最小值，最小值
\par
丿0）一2，最小值点为1.
\end{ocrdraft}
\end{bankitem}
